%% --------------------------------------------------------
\section{CASSCF --- багатоконфігураційний метод}
%% --------------------------------------------------------

%% --------------------------------------------------------
\subsection{Ідея Complete Active Space Self-Consistent Field}
%% --------------------------------------------------------

CASSCF (Complete Active Space SCF) --- багатоконфігураційний метод, який поєднує:
\begin{itemize}
    \item \textbf{Активний простір} --- набір орбіталей, де проводиться Full CI
    \item \textbf{SCF оптимізацію} --- орбіталі та CI коефіцієнти оптимізуються одночасно
\end{itemize}

Орбіталі розділяються на три групи:
\begin{enumerate}
    \item \textbf{Core (остов)} --- завжди подвійно заповнені, не беруть участь у кореляції
    \item \textbf{Active (активні)} --- електрони розподілені по всіх можливих способах (Full CI)
    \item \textbf{Virtual (віртуальні)} --- завжди порожні
\end{enumerate}

\begin{equation}
    |\Psi_{CASSCF}\rangle = \sum_I c_I |\Phi_I^{\text{active}}\rangle \otimes |\text{core}\rangle
\end{equation}

%% --------------------------------------------------------
\subsection{Позначення CASSCF(n,m)}
%% --------------------------------------------------------

CASSCF(n,m) означає:
\begin{itemize}
    \item $n$ --- кількість електронів в активному просторі
    \item $m$ --- кількість орбіталей в активному просторі
\end{itemize}

\textbf{Приклади:}
\begin{itemize}
    \item \textbf{Be}: CASSCF(4,8) --- 4 валентні електрони у 8 орбіталях (2s, 2p, 3s, 3p)
    \item \textbf{C}: CASSCF(4,4) --- 4 валентні електрони у 4 орбіталях (2s, 2p)
    \item \textbf{Cr}: CASSCF(6,5) --- 6 електронів у 5 d-орбіталях
\end{itemize}

%% --------------------------------------------------------
\subsection{CASSCF розрахунок атома Берилію}
%% --------------------------------------------------------

Берилій має близькі за енергією конфігурації 2s² та 2p², що робить його класичним прикладом для CASSCF:

\inputcode{CASSCFBe.py}

%% --------------------------------------------------------
\subsection{CASSCF для перехідних металів}
%% --------------------------------------------------------


Для перехідних металів d-орбіталі часто близькі за енергією, потребуючи багатоконфігураційного підходу:

\inputcode{CASSCFTransitionMetals.py}


%% --------------------------------------------------------
\subsection{Вибір активного простору}
%% --------------------------------------------------------

Правильний вибір активного простору критичний для CASSCF:

\begin{table}[h]
    \centering
    \caption{Рекомендації по вибору активного простору}
    \label{tab:casscf_active_space}
    \small
    \begin{tabular}{lll}
        \hline
        \textbf{Система} & \textbf{Активний простір} & \textbf{Коментар}          \\
        \hline
        Be               & (4,8): 2s,2p,3s,3p        & Враховує $2s^2 \leftrightarrow 2p^2$           \\
        C                & (4,4): 2s,2p              & Мінімальний валентний      \\
        O                & (6,6): 2s,2p + корел.     & З кореляційними орбіталями \\
        3d метали        & (n,5): 3d                 & Тільки d-орбіталі          \\
        3d метали        & (n+2,6): 3d,4s            & d + s орбіталі             \\
        Лантаноїди       & (n,7): 4f                 & f-орбіталі                 \\
        \hline
    \end{tabular}
\end{table}

\textbf{Принципи вибору:}
\begin{enumerate}
    \item Включити орбіталі, близькі за енергією
    \item Включити орбіталі з значною хімічною активністю
    \item Більший простір $\rightarrow$ точніше, але експоненційно повільніше
    \item Перевірити заселеності природних орбіталей
    \item Якщо заселеності $\approx2$ або $\approx0$, можна виключити з активного простору
\end{enumerate}


%% --------------------------------------------------------
\subsection{State-averaged CASSCF}
%% --------------------------------------------------------

Для розрахунку декількох станів одночасно (наприклад, різних мультиплетів):

\inputcode{StateAvaragedCASSCF.py}

%% --------------------------------------------------------
\subsection{CASPT2 --- динамічна кореляція поверх CASSCF}
%% --------------------------------------------------------

CASSCF враховує тільки статичну кореляцію. Для повної точності потрібно додати динамічну кореляцію через CASPT2:

\inputcode{CASPT2.py}

%% --------------------------------------------------------
\subsection{Переваги та недоліки CASSCF}
%% --------------------------------------------------------

\paragraph{Переваги:}
\begin{itemize}
    \item Правильно описує статичну кореляцію
    \item Може розрахувати декілька станів одночасно
    \item Варіаційний метод
    \item Підходить для розриву зв'язків, збуджених станів
    \item Дає природні орбіталі та заселеності
\end{itemize}

\paragraph{Недоліки:}
\begin{itemize}
    \item Вибір активного простору не завжди очевидний
    \item Експоненційне масштабування з розміром активного простору
    \item Не враховує динамічну кореляцію (потрібен CASPT2)
    \item Може бути проблеми з конвергенцією
    \item Потребує досвіду для правильного використання
\end{itemize}

\paragraph{Коли використовувати CASSCF:}
\begin{itemize}
    \item T1 діагностика > 0.05 (сильна багатоконфігураційність)
    \item Перехідні метали з близькими d-орбіталями
    \item Збуджені стани атомів
    \item Діелектронні системи
    \item Відкрито-оболонкові синглети
\end{itemize}
